\section{Conclusion}
\label{sec:conclusion}

This article evaluated adaptive and cost-aware data quality monitoring for evolving batch pipelines. The monitor combines numerical adaptation, four-way batch labelling, historical admission of LLM-proposed rules, category-admission guards and risk-based check selection. Its evaluation addresses the three research questions using prespecified criteria, compares components with calibrated alternatives and configured tools, and supplements controlled injection with dirty/clean benchmarks and chronological replay.

The results identify where each component contributes. Four-way macro-F1 reaches 0.947--0.966, but fixed thresholds behind the same gated drift signal perform identically on Chicago and show no detected disadvantage on Retail. Numerical adaptation contributes most clearly on BTS, reducing drift-induced false alarms, while the small Chicago benefit comes from category policy. Retail does not meet the lower-FPR hypothesis. Persistent contamination and invalid-category recurrence demonstrate that adaptive expectations can also weaken error detection.

Domain and statistical rules are complementary. Fixed statistical rules with admitted LLM rules reach semantic recall of 0.996--1.00 on the selected injected violations with at most a 0.01 pp primary-model FPR increase. This result does not describe the adaptive combination, which remains vulnerable to recurring invalid codes. A separate per-value policy admits all twelve valid novel values and rejects eleven of twelve invalid values. Tax further shows that cross-column relationships matter: learned dependencies raise cell recall within checkable columns from about 0.20 to 0.91, with the primary LLM's structure matching rather than exceeding that record recall.

Risk-based selection is useful in eleven of eighteen degradation/incident cases by the specified mean-recall criterion, but saves only about 2--11\% of checks. Profiling and drift work dominate total cost, and selection adds no material timing benefit after profile reuse in the measured workloads. The chronological Chicago replay labels the documented July 2018 procedure change as drift without degradation, while also revealing substantial alert burden and incomplete human ground truth.

The limitations identify four remaining tasks within the existing research scope. Human adjudication of replay alarms and expert-authored comparison rules would strengthen semantic interpretation. Pipelines with logged natural incidents would test the injected definition of degradation. Improving the efficiency of profiling and drift calculation would address the dominant observed cost. Finally, evaluating cross-column rules within the adaptive chronological loop would extend the existing benchmark analysis to the monitored setting. These are follow-up validation needs, not completed results.

The central implication is that adaptation and validation must remain distinguishable. A monitor should follow legitimate changes while preserving independently justified constraints on valid data, and its detection and cost benefits should be assessed for the exact configuration used.

\section*{Data and Code Availability}
The three source datasets and two benchmark tables are public~\cite{BTS2024,Chen2019,Chicago2026,Mahdavi2019}. The research artifact contains the monitoring implementation, experiment harness, frozen configurations, cached responses, run manifests and result tables. These materials are intended for public release with the final article. A public archival repository identifier is not specified in this version.
